\chapter{Senyawa Karbonil: Adisi Nukleofilik, Asetal, dan Proteksi}\label{ch:b1:acetals-protection}

Larutkan glukosa dalam air, maka hampir tidak ada yang tetap berupa
aldehida rantai terbuka seperti yang digambar di buku pelajaran: molekul
itu melengkung, gugus hidroksilnya sendiri beradisi pada gugus
aldehidanya, dan sebuah cincin beranggota enam menutup. Adisi alkohol
yang sama pada gugus karbonil, yang diulang sekali lagi dengan \omterm{def:b1:elementary-steps:catalyst}{katalis}
asam, menghasilkan \omterm{def:b1:acetals-protection:hemiacetal}{asetal}, senyawa yang stabil terhadap basa dan terhadap
\omterm{def:b1:electronic-effects:nucleophile}{nukleofil} yang paling reaktif. Kestabilan itu adalah sebuah alat: aldehida
yang akan dirusak oleh \omterm{def:b1:organomagnesium:organometallic}{pereaksi Grignard} dapat disembunyikan sebagai
\omterm{def:b1:acetals-protection:hemiacetal}{asetal}, dibawa melewati reaksi itu, dan diperoleh kembali pada akhirnya.
Bab ini mempelajari gugus karbonil sebagai \omterm{def:b1:electronic-effects:nucleophile}{elektrofil}, adisi air dan
alkohol, serta strategi \omterm{def:b1:acetals-protection:protecting-group}{proteksi}.

\begin{recall}
Adisi Grignard pada senyawa karbonil (\cref{ch:b1:organomagnesium});
\omterm{def:b1:extent-q-and-k:quotient}{kuosien reaksi} $Q$ dan evolusi sistem selama $Q \neq K$
(\cref{ch:b1:extent-q-and-k}); \omterm{def:b1:electronic-effects:curly-arrow}{panah lengkung}, \omterm{def:b1:electronic-effects:nucleophile}{nukleofil}, dan \omterm{def:b1:electronic-effects:nucleophile}{elektrofil}
(\cref{ch:b1:electronic-effects}). Jilid sekolah (kelas 12) telah
memperkenalkan gagasan melindungi sebuah gugus selama sintesis.
\end{recall}

\section{Gugus karbonil sebagai elektrofil}

Ikatan C=O terpolarisasi ke arah oksigen (\cref{ch:b1:periodicity}), dan
\omterm{def:b1:lewis-resonance-vsepr:resonance}{struktur resonansi} \ce{C+-O-} menempatkan muatan positif penuh pada
karbon: karbon itu elektrofilik, diserang oleh \omterm{def:b1:electronic-effects:nucleophile}{nukleofil} tegak lurus
terhadap bidang gugus itu. \omterm{def:b1:electronic-effects:nucleophile}{Nukleofil} kuat (\omterm{def:b1:organomagnesium:organometallic}{pereaksi Grignard}, hidrida,
sianida) beradisi secara langsung. \omterm{def:b1:electronic-effects:nucleophile}{Nukleofil} lemah --- air, alkohol ---
memerlukan bantuan.

\begin{definition}[Aktivasi elektrofilik]\label{def:b1:acetals-protection:activation}
\emph{Aktivasi elektrofilik}\index{aktivasi elektrofilik} gugus karbonil
adalah protonasinya (atau koordinasinya dengan \omterm{def:b1:lewis-resonance-vsepr:lewis-acid}{asam Lewis}) pada oksigen.
Kation \ce{C=OH+}, yang muatan positifnya ikut dimiliki karbon melalui
\omterm{def:b1:lewis-resonance-vsepr:resonance}{struktur resonansi} \ce{C+-OH}, diserang bahkan oleh \omterm{def:b1:electronic-effects:nucleophile}{nukleofil} yang
lemah.
\end{definition}

\begin{omfigure}
\setchemfig{atom sep=2.1em}
\schemestart
\chemfig{C(-[:150]R)(-[:210]H)=O}
\arrow{->[\ce{H+}]}
\chemfig{C(-[:150]R)(-[:210]H)=O^{+}-H}
\arrow{<->}
\chemfig{C^{+}(-[:150]R)(-[:210]H)-O-H}
\schemestop

{\small Aktivasi aldehida oleh asam: protonasi oksigen menghasilkan
kation yang \omterm{def:b1:lewis-resonance-vsepr:resonance}{struktur resonansinya} membawa muatan pada karbon.}
\end{omfigure}

\section{Hidrat dan hemiasetal}

Air beradisi secara reversibel pada aldehida dan keton, \ce{R2C=O + H2O <=>
R2C(OH)2}; hidratnya biasanya minor untuk keton dan dapat mayor untuk
metanal. Alkohol beradisi dengan cara yang sama.

\begin{definition}[Hemiasetal dan asetal]\label{def:b1:acetals-protection:hemiacetal}
\emph{Hemiasetal}\index{hemiasetal} adalah senyawa yang satu karbonnya
membawa sebuah gugus OH sekaligus sebuah gugus OR, yang terbentuk melalui
adisi alkohol pada gugus karbonil. \emph{Asetal}\index{asetal} adalah
senyawa yang satu karbonnya membawa dua gugus OR, yang terbentuk dari
hemiasetal dan molekul alkohol kedua, disertai lepasnya air.
\end{definition}

\omterm{def:b1:acetals-protection:hemiacetal}{Hemiasetal} rantai terbuka biasanya minor dalam kesetimbangannya dengan
aldehida dan alkohol; jika OH dan C=O berada dalam molekul yang sama dan
sebuah cincin beranggota lima atau enam dapat menutup, \omterm{def:b1:acetals-protection:hemiacetal}{hemiasetal} siklik
itulah yang predominan. Demikianlah halnya pada glukosa, dan pada
5-hidroksipentanal yang lebih sederhana.

\begin{omfigure}
\setchemfig{atom sep=2em}
\schemestart
\chemfig{HO-[:30]-[:-30]-[:30]-[:-30]-[:30]C(=[:90]O)-[:-30]H}
\arrow{<=>}
\chemfig{*6(O-(-[:-30]OH)-----)}
\schemestop

\medskip

{\small 5-Hidroksipentanal dan \omterm{def:b1:acetals-protection:hemiacetal}{hemiasetal} siklisnya, sebuah cincin
beranggota enam yang mengandung oksigen. OH karbon 5 telah beradisi pada
aldehida karbon 1; bentuk cincinnya predominan, seperti pada glukosa.}
\end{omfigure}

\section{Asetal}

\begin{proposition}[Mekanisme pembentukan asetal]\label{prop:b1:acetals-protection:acetal-mechanism}
Dalam asam, aldehida atau keton dan alkohol menghasilkan \omterm{def:b1:acetals-protection:hemiacetal}{hemiasetal}, lalu
\omterm{def:b1:acetals-protection:hemiacetal}{asetal}, melalui tahap-tahap: protonasi C=O; adisi alkohol; lepasnya proton
(\omterm{def:b1:acetals-protection:hemiacetal}{hemiasetal}); protonasi OH; lepasnya air, yang menghasilkan \omterm{def:b1:electronic-effects:intermediates}{karbokation}
yang distabilkan oleh gugus OR yang tersisa (ion oksokarbenium); adisi
alkohol kedua; lepasnya proton. Setiap tahap bersifat reversibel.
\end{proposition}

\begin{proof}
Setiap tahap adalah perpindahan proton asam--basa atau adisi (atau
lepasnya) sebuah \omterm{def:b1:electronic-effects:nucleophile}{nukleofil} pada (atau dari) karbon yang elektrofilik,
seperti yang diuraikan dalam \cref{ch:b1:electronic-effects}. Lepasnya air
dimungkinkan karena kation yang terbentuk, $\mathrm{R_2C^+{-}OR}$, mempunyai \omterm{def:b1:lewis-resonance-vsepr:resonance}{struktur resonansi} $\mathrm{R_2C{=}O^+R}$
yang setiap atomnya memiliki oktet (\cref{prop:b1:electronic-effects:carbocation-order}).
Asam dipakai dalam protonasi dan dibentuk kembali dalam deprotonasi:
asam itu adalah \omterm{def:b1:elementary-steps:catalyst}{katalis}.
\end{proof}

\begin{omfigure}
\setchemfig{atom sep=1.9em}
\schemestart
\chemfig{C(-[:150]R)(-[:210]H)=O}
\arrow{->[\ce{H+}][]}
\chemfig{C^{+}(-[:150]R)(-[:210]H)-OH}
\arrow{->[$\mathrm{R'OH}$][\ce{-H+}]}
\chemfig{C(-[:150]R)(-[:210]H)(-[2]OR')-OH}
\schemestop

\medskip
\schemestart
\chemfig{C(-[:150]R)(-[:210]H)(-[2]OR')-OH}
\arrow{->[\ce{H+}][\ce{-H2O}]}
\chemfig{C^{+}(-[:150]R)(-[:210]H)-OR'}
\arrow{->[$\mathrm{R'OH}$][\ce{-H+}]}
\chemfig{C(-[:150]R)(-[:210]H)(-[2]OR')-OR'}
\schemestop
\par\vspace{6pt}

{\small Pembentukan \omterm{def:b1:acetals-protection:hemiacetal}{asetal}, dalam dua baris: mula-mula \omterm{def:b1:acetals-protection:hemiacetal}{hemiasetal}
(aktivasi, adisi alkohol, lepasnya proton), lalu \omterm{def:b1:acetals-protection:hemiacetal}{asetal} (protonasi OH,
lepasnya air menjadi kation yang stabil, adisi alkohol kedua, lepasnya
proton).}
\end{omfigure}

Reaksi keseluruhan, $\mathrm{RCHO} + 2\,\mathrm{R'OH} \rightleftharpoons \mathrm{RCH(OR')_2} + \ce{H2O}$, mempunyai
tetapan kesetimbangan yang sedang, dan air adalah salah satu produknya.

\begin{proposition}[Mendorong asetalisasi]\label{prop:b1:acetals-protection:dean-stark}
Jika air yang terbentuk dikeluarkan secara terus-menerus, asetalisasi
berlanjut sampai senyawa karbonil (atau alkoholnya) habis.
\end{proposition}

\begin{proof}
\omterm{def:b1:extent-q-and-k:quotient}{Kuosien reaksi} $Q = [\text{asetal}][\ce{H2O}]/([\text{aldehida}]
[\mathrm{R'OH}]^2)$ tetap di bawah $K$ selama air dikeluarkan: menurut
\cref{ch:b1:extent-q-and-k}, reaksi terus bergerak maju, berapa pun nilai
$K$. Dengan sebuah diol (etana-1,2-diol) \omterm{def:b1:acetals-protection:hemiacetal}{asetalnya} siklik, dan
kesetimbangannya juga lebih menguntungkan, karena satu molekul diol
menggantikan dua molekul alkohol.
\end{proof}

\begin{omfigure}
\begin{tikzpicture}[font=\small]
  \pic at (0,0) {heatingmantle};
  \pic at (0,0) {rbflask={0.4}{omDef!14}};
  \pic (d) at (0,3.0) {deanstark={0.35}{omDef!35}{orange!15}};
  \pic at (0,6.4) {condenser};
  \draw[->] (2.4,4.0) node[right, align=left] {penampung: air (lapisan bawah)\\tidak kembali lagi; toluena\\meluap kembali ke labu} -- (0.8,3.8);
  \draw[->] (-2.2,1.0) node[left, align=right] {aldehida, diol, toluena,\\katalis asam, mendidih} -- (-0.7,0.9);
  \draw[->] (1.8,8.0) node[right] {kondensor} -- (0.4,7.8);
\end{tikzpicture}

{\small Perangkap Dean--Stark. Uap toluena dan air mengembun dan jatuh ke
dalam penampung bergaris skala; air, yang lebih rapat dan tidak \omterm{def:b1:intermolecular-forces-solvents:miscibility}{dapat
bercampur} dengan toluena, tenggelam dan tinggal di sana, sedangkan
toluena meluap kembali ke dalam labu. Volume air yang terkumpul mengukur
\omterm{def:b1:extent-q-and-k:extent}{kemajuan reaksi}.}
\end{omfigure}

\begin{proposition}[Kestabilan asetal]\label{prop:b1:acetals-protection:acetal-stability}
\omterm{def:b1:acetals-protection:hemiacetal}{Asetal} stabil dalam medium netral dan basa serta terhadap \omterm{def:b1:electronic-effects:nucleophile}{nukleofil},
\omterm{def:b1:organomagnesium:organometallic}{pereaksi Grignard}, dan pereaksi hidrida; asam berair menghidrolisisnya
kembali menjadi senyawa karbonil dan alkohol.
\end{proposition}

\begin{proof}
Karbon \omterm{def:b1:acetals-protection:hemiacetal}{asetal} tidak membawa \omterm{def:b1:electronic-effects:nucleophile}{gugus pergi} yang dapat diusir oleh basa atau
\omterm{def:b1:electronic-effects:nucleophile}{nukleofil}: \ce{RO-}, sebuah \omterm{def:b1:predominant-reaction:strong-acid}{basa kuat}, tidak pergi, dan tidak ada lagi
C=O yang dapat diserang. Dalam asam berair, mekanisme di atas berjalan
mundur: air yang sangat berlebih membuat $Q < K$ untuk reaksi balik.
\end{proof}

\section{Proteksi dan deproteksi}

\begin{definition}[Gugus pelindung]\label{def:b1:acetals-protection:protecting-group}
\emph{Gugus pelindung}\index{gugus pelindung} adalah gugus yang dimasukkan
ke dalam molekul untuk menutupi sebuah fungsi untuk sementara, agar fungsi
itu tidak bereaksi dalam tahap yang ditujukan pada bagian lain molekul.
Pemasukannya adalah \emph{proteksi}\index{proteksi}, pelepasannya adalah
\emph{deproteksi}\index{deproteksi}.
\end{definition}

\begin{method}[Proteksi, reaksi, deproteksi]\label{met:b1:acetals-protection:protect}
\begin{enumerate}
  \item Kenali fungsi yang akan bereaksi, atau yang akan merusak
        pereaksinya, dalam tahap yang direncanakan.
  \item Pilihlah \omterm{def:b1:acetals-protection:protecting-group}{gugus pelindung} yang stabil dalam tahap itu dan dapat
        dilepas pada kondisi yang dapat ditoleransi bagian molekul
        lainnya (untuk aldehida atau keton yang menghadapi basa,
        \omterm{def:b1:electronic-effects:nucleophile}{nukleofil}, atau hidrida: \omterm{def:b1:acetals-protection:hemiacetal}{asetal} siklik).
  \item Lakukan \omterm{def:b1:acetals-protection:protecting-group}{proteksi}; jalankan tahap yang direncanakan; lakukan
        \omterm{def:b1:acetals-protection:protecting-group}{deproteksi}.
  \item Hitunglah ongkosnya: dua tahap tambahan, masing-masing dengan
        \omterm{def:b1:extent-q-and-k:fractional-extent}{rendemennya}.
\end{enumerate}
\end{method}

\begin{example}[Keton yang menghalangi]\label{ex:b1:acetals-protection:ketoester}
Untuk mengadisikan \omterm{def:b1:organomagnesium:organometallic}{pereaksi Grignard} pada gugus aldehida sebuah molekul
yang juga membawa keton, tidak ada pilihan \omterm{def:b1:acetals-protection:hemiacetal}{asetal} yang sederhana (aldehida
lebih mudah membentuk \omterm{def:b1:acetals-protection:hemiacetal}{asetal} daripada keton, sehingga gugus yang salah
yang akan terlindungi); urutan tahap atau pereaksinya harus diubah.
\omterm{def:b1:acetals-protection:protecting-group}{Proteksi} adalah sebuah strategi, bukan refleks.
\end{example}

\section{Latihan}

\begin{exercise}[$\star$]\label{exo:b1:acetals-protection:1}
Tentukan karbon \omterm{def:b1:acetals-protection:hemiacetal}{hemiasetal}, \omterm{def:b1:acetals-protection:hemiacetal}{asetal}, hidrat, eter, atau alkohol dalam:
\ce{CH3CH(OH)OCH3}, \ce{CH3CH(OCH3)2}, \ce{CH3CH(OH)2}, \ce{CH3OCH3},
\ce{(CH3)2C(OCH2CH3)2}.
\end{exercise}

\begin{exercise}[$\star$]\label{exo:b1:acetals-protection:2}
Tuliskan persamaan keseluruhan pembentukan \omterm{def:b1:acetals-protection:hemiacetal}{asetal} propanal dengan
metanol, dan \omterm{def:b1:acetals-protection:hemiacetal}{asetal} siklik propanon dengan etana-1,2-diol.
\end{exercise}

\begin{exercise}[$\star$]\label{exo:b1:acetals-protection:3}
Pereaksi manakah di antara berikut ini yang membiarkan \omterm{def:b1:acetals-protection:hemiacetal}{asetal} tidak
berubah: larutan natrium hidroksida, metilmagnesium bromida, asam sulfat
encer dalam air, natrium borohidrida?
\end{exercise}

\begin{exercise}[$\star$]\label{exo:b1:acetals-protection:4}
Gambarlah \omterm{def:b1:acetals-protection:hemiacetal}{hemiasetal} siklik yang dibentuk oleh 4-hidroksibutanal.
Berapakah ukuran cincinnya?
\end{exercise}

\begin{exercise}[$\star\star$]\label{exo:b1:acetals-protection:5}
Tuliskan mekanisme lengkap pembentukan dimetil \omterm{def:b1:acetals-protection:hemiacetal}{asetal} etanal yang
dikatalisis asam, dengan \omterm{def:b1:electronic-effects:curly-arrow}{panah lengkung}, dan tunjukkan bahwa asamnya
terbentuk kembali.
\end{exercise}

\begin{exercise}[$\star\star$]\label{exo:b1:acetals-protection:6}
Tuliskan mekanisme hidrolisis asam 2,2-dimetoksipropana. Mengapa dipakai
air dalam jumlah besar?
\end{exercise}

\begin{exercise}[$\star\star$]\label{exo:b1:acetals-protection:7}
Asetalisasi \qty{0.250}{mol} sebuah keton diikuti dengan perangkap
Dean--Stark. Berapa volume air yang seharusnya terkumpul pada konversi
penuh (massa jenis sekitar \qty{1.0}{g/mL})? Sesudah satu jam,
\qty{3.2}{mL} telah terkumpul: berapakah konversinya?
\end{exercise}

\begin{exercise}[$\star\star$]\label{exo:b1:acetals-protection:8}
Jelaskan mengapa \omterm{def:b1:acetals-protection:hemiacetal}{asetal} siklik dengan etana-1,2-diol terbentuk lebih
sempurna daripada \omterm{def:b1:acetals-protection:hemiacetal}{asetal} dengan dua molekul metanol, untuk senyawa
karbonil yang sama.
\end{exercise}

\begin{exercise}[$\star\star$]\label{exo:b1:acetals-protection:9}
Mengapa \omterm{def:b1:elementary-steps:catalyst}{katalis} asam harus dihilangkan (atau dinetralkan) sebelum senyawa
yang terlindungi direaksikan dengan \omterm{def:b1:organomagnesium:organometallic}{pereaksi Grignard}?
\end{exercise}

\begin{exercise}[$\star\star\star$]\label{exo:b1:acetals-protection:10}
Usulkan sintesis 5-hidroksi-5-metilheksan-2-on,
\ce{(CH3)2C(OH)CH2CH2COCH3}, dari 4-klorobutan-2-on dan propanon, dengan
memakai sebuah \omterm{def:b1:acetals-protection:protecting-group}{proteksi}. Tentukan setiap tahap, pereaksinya, dan peran
masing-masing.
\end{exercise}

\begin{exercise}[$\star\star\star$]\label{exo:b1:acetals-protection:11}
Tunjukkan dengan \omterm{def:b1:extent-q-and-k:quotient}{kuosien reaksi} bahwa metanol yang sangat berlebih juga
mendorong asetalisasi sebuah aldehida, dan bandingkan metode ini dengan
pengeluaran air.
\end{exercise}

\begin{exercise}[$\star\star\star$]\label{exo:b1:acetals-protection:12}
Jelaskan mengapa \omterm{def:b1:acetals-protection:hemiacetal}{hemiasetal} siklik glukosa diserang, dalam metanol yang
asam, menghasilkan metil \omterm{def:b1:acetals-protection:hemiacetal}{asetal} (sebuah metil glikosida), sedangkan gugus
OH glukosa yang lain tidak diubah menjadi eter.
\end{exercise}

\section{Soal: Pereaksi Grignard yang Membawa Aldehida}

\begin{problem}[{Soal akhir pekan --- mengapa 4-bromobutanal tidak dapat
menghasilkan pereaksi Grignard, proteksinya sebagai asetal siklik dengan
perangkap Dean--Stark, adisi Grignard pada propanon, deproteksinya, dan
volume air yang terkumpul}]
\label{pb:b1:acetals-protection:1}
Seorang kimiawan memerlukan 5-hidroksi-5-metilheksanal,
\ce{(CH3)2C(OH)CH2CH2CH2CHO}, dan berencana membuatnya dari
4-bromobutanal, \ce{BrCH2CH2CH2CHO}, dan propanon. Ia memulai dari
\qty{15.09}{g} 4-bromobutanal. Massa molar (\unit{g/mol}): \ce{C4H7BrO}
150.9, \ce{C2H6O2} 62.0, \ce{C6H11BrO2} 194.9, \ce{C3H6O} 58.0,
\ce{C9H18O3} 174.0, \ce{C7H14O2} 130.0, \ce{H2O} 18.0; massa jenis air
\qty{0.997}{g/mL}.
% ledger: rho:water

\medskip
\noindent\textbf{Bagian I --- Persoalannya.}
\begin{enumerate}
  \item Tuliskan \omterm{def:b1:organomagnesium:organometallic}{pereaksi Grignard} yang akan dihasilkan 4-bromobutanal.
  \item Mengapa pereaksi itu tidak dapat ada? Tuliskan reaksi yang
        merusaknya.
  \item Fungsi manakah yang harus diproteksi, dan terhadap apa?
  \item Mengapa \omterm{def:b1:acetals-protection:hemiacetal}{asetal} cocok?
  \item Mengapa \omterm{def:b1:acetals-protection:hemiacetal}{asetal} siklik (dengan etana-1,2-diol) lebih disukai?
  \item Hitunglah jumlah 4-bromobutanal.
  \item Hitunglah massa diol untuk kelebihan \qty{20}{\percent}.
\end{enumerate}

\medskip
\noindent\textbf{Bagian II --- \omterm{def:b1:acetals-protection:protecting-group}{Proteksi}.}
\begin{enumerate}[resume]
  \item Tuliskan persamaan asetalisasi dengan etana-1,2-diol.
  \item Apa peran \omterm{def:b1:elementary-steps:catalyst}{katalis} asam (asam 4-metilbenzenasulfonat)?
  \item Uraikan cara kerja perangkap Dean--Stark.
  \item Jelaskan dengan $Q$ dan $K$ mengapa pengeluaran air mendorong
        reaksinya.
  \item Mengapa toluena, dan bukan air, yang dikembalikan ke labu?
  \item Bagaimana kimiawan itu mengetahui bahwa reaksinya sudah selesai?
  \item Hitunglah massa bromida terlindungi yang diharapkan pada konversi
        penuh.
\end{enumerate}

\medskip
\noindent\textbf{Bagian III --- Tahap Grignard.}
\begin{enumerate}[resume]
  \item Tuliskan pembentukan \omterm{def:b1:organomagnesium:organometallic}{pereaksi Grignard} dari bromida yang
        terlindungi itu.
  \item Tuliskan adisinya pada propanon dan \omterm{def:b1:alcohol-activation:alkoxide}{alkoksida} yang terbentuk.
  \item Mengapa pengolahan akhir tahap ini harus bersifat agak basa atau
        netral (misalnya amonium klorida berair) dan bukan asam?
  \item Hitunglah massa alkohol terlindungi untuk \omterm{def:b1:extent-q-and-k:fractional-extent}{rendemen}
        \qty{70}{\percent} pada tahap ini.
\end{enumerate}

\medskip
\noindent\textbf{Bagian IV --- \omterm{def:b1:acetals-protection:protecting-group}{Deproteksi}.}
\begin{enumerate}[resume]
  \item Tuliskan \omterm{def:b1:acetals-protection:protecting-group}{deproteksinya} dalam asam berair.
  \item Mengapa alkohol tersier itu bertahan dalam asam berair yang lunak?
  \item Produknya dapat menutup cincin beranggota enam dengan
        mengadisikan OH-nya pada aldehida. Gambarlah \omterm{def:b1:acetals-protection:hemiacetal}{hemiasetal} siklik itu.
  \item Hitunglah massa hidroksi aldehida pada \omterm{def:b1:extent-q-and-k:fractional-extent}{rendemen} \omterm{def:b1:acetals-protection:protecting-group}{deproteksi}
        \qty{90}{\percent}.
  \item Berapa \omterm{def:b1:extent-q-and-k:fractional-extent}{rendemen} keseluruhan dari 4-bromobutanal yang dicapai
        kimiawan itu?
  \item Hitunglah massa air yang terbentuk pada tahap \omterm{def:b1:acetals-protection:protecting-group}{proteksi} pada
        konversi penuh.
  \item Nyatakan volume air yang terkumpul dalam perangkap Dean--Stark
        pada konversi penuh.
\end{enumerate}
\end{problem}
